\chapter*{Il \textit{Devī Sūkta }e la dimensione tantrica}
\addcontentsline{toc}{chapter}{Il \textit{Devī Sūkta }e la dimensione tantrica}

Parlare di un “aspetto tantrico” del \textit{Devī Sūkta} richiede cautela. L’inno appartiene pienamente al \textit{Ṛgveda} e precede di secoli la formazione storica delle tradizioni tantriche come sistemi riconoscibili. Tuttavia, alcune strutture di esperienza e di linguaggio che il tantrismo svilupperà in forma esplicita sono già presenti, in modo germinale e non sistematico, in questo testo vedico. Non si tratta quindi di leggere il \textit{Devī Sūkta} \emph{come} un testo tantrico, ma di riconoscere in esso una sensibilità che il tantra farà propria.

Il primo elemento decisivo è la voce in prima persona. Nel \textit{Devī Sūkta} non si descrive la Dea dall’esterno, né la si invoca come oggetto di culto: la Dea parla. Questo spostamento del punto di vista è radicale. La realtà ultima non è qualcosa da raggiungere, ma qualcosa che prende parola dall’interno dell’esperienza. Il tantrismo erediterà proprio questa postura: non una trascendenza distante, ma una presenza che si manifesta come soggetto.

In secondo luogo, la Dea non viene collocata in un ambito separato dal mondo. Lei afferma di muoversi con gli dèi, di abitare gli atti più elementari della vita — mangiare, respirare, ascoltare — e di pervadere cielo e terra. Questa assenza di distinzione tra sacro e profano è una delle consonanze più profonde con la visione tantrica. Ciò che è vitale non è da purificare per diventare sacro: è già il luogo della manifestazione del sacro, se riconosciuto.

Un tratto particolarmente affine alla sensibilità tantrica è il modo in cui il \textit{Devī Sūkta }tratta il potere. La Dea non è soltanto origine o sostegno: ella conferisce efficacia. Rende qualcuno \textit{ugra}, lo fa \textit{ṛṣi}, \textit{brahmāṇa}, \textit{sumedhā}. Conoscenza, parola autorevole e forza non sono separabili. Nel tantrismo, questa unità diventerà centrale: la conoscenza autentica non è neutra, ma trasformativa; produce effetti, modifica la postura dell’essere nel mondo.

Anche il rapporto con il desiderio e il conflitto mostra una prossimità significativa. Quando la Dea tende l’arco per \textit{Rudra} o prepara il \textit{samada}, il confronto, non introduce una violenza caotica, ma una forza che ristabilisce l’ordine. Il tantra non elimina le energie intense della vita, ma le assume come strumenti di trasformazione. In questo senso, il \textit{Devī Sūkta} offre una visione del potere che non è moraleggiante, ma funzionale all’equilibrio del cosmo.

Un altro punto di contatto fondamentale è la centralità della parola. La Dea è \textit{Vāk}, ma non solo come linguaggio articolato: è parola efficace, performativa, che crea, sostiene e trasforma. Il tantrismo svilupperà questa intuizione nella dottrina del \textit{mantra}, ma nel \textit{Devī Sūkta} la parola non è ancora tecnica rituale: è potenza originaria che si esprime prima di ogni formalizzazione. Il \textit{mantra} non viene ancora insegnato; viene mostrata la sua sorgente.

Infine, il \textit{Devī Sūkta }presenta una visione del divino che non esclude l’ignoranza umana, ma la include. Molti “dimorano vicino” alla Dea senza riconoscerla. Questa constatazione non è un giudizio, ma un dato di esperienza. La distanza non è spaziale, bensì percettiva. Anche questo sarà un tema centrale del tantrismo: la liberazione non consiste nel raggiungere un luogo altro, ma nel mutare il modo di vedere ciò che è già presente.

In questo senso, il \textit{Devī Sūkta} può essere letto come una soglia. Non è ancora tantra, ma rende possibile una visione in cui il corpo, la parola, il potere, il desiderio e il cosmo non sono realtà da superare, bensì da attraversare consapevolmente. Il tantrismo non nasce contro il \textit{Veda}: nasce anche da intuizioni come questa, lasciate aperte, non sistematizzate, ma potenti nella loro semplicità.

